\subsection{Kronecker Products and Vectorization}

Throughout, $\mathbb{F}$ denotes a field.

\begin{definition} \label{def:kronecker-product} \hlt{Kronecker Product}\\
Let $m,n,p,q\in\mathbb{Z}_{>0}$, $A=[a_{ij}]\in\operatorname{Mat}(m,n,\mathbb{F})$, and $B\in\operatorname{Mat}(p,q,\mathbb{F})$.\\
The \hlt{Kronecker product}, or \hlt{matrix tensor product}, is the block matrix $A\otimes B\defeq[a_{ij}B]_{i\in J_m,j\in J_n}\in\operatorname{Mat}(mp,nq,\mathbb{F})$.
\end{definition}

\begin{proposition} \label{prop:basic-kronecker-identities} \hlt{Basic Kronecker Identities}\\
Let $m,n,p,q,r,s\in\mathbb{Z}_{>0}$, $c\in\mathbb{F}$, and $A,A'\in\operatorname{Mat}(m,n,\mathbb{F})$.\\
Let $B,B'\in\operatorname{Mat}(p,q,\mathbb{F})$ and $C\in\operatorname{Mat}(r,s,\mathbb{F})$.
\begin{enumerate}[label=(\roman*)]
\setlength{\itemsep}{0pt}
\item $c(A\otimes B)=(cA)\otimes B=A\otimes(cB)$;
\item $(A+A')\otimes B=A\otimes B+A'\otimes B$ and $A\otimes(B+B')=A\otimes B+A\otimes B'$;
\item $(A\otimes B)\otimes C=A\otimes(B\otimes C)$;
\item $(A\otimes B)^t=A^t\otimes B^t$;
\item $I_m\otimes I_p=I_{mp}$;
\item If $\mathbb{F}=\mathbb{C}$, then $\overline{A\otimes B}=\overline{A}\otimes\overline{B}$ and $(A\otimes B)^*=A^*\otimes B^*$.
\end{enumerate}
\begin{proof}
Write $A=[a_{ij}]$ and $B=[b_{uv}]$.
\begin{enumerate}[label=(\roman*)]
\setlength{\itemsep}{0pt}
\item The corresponding blocks in the three matrices are $c(a_{ij}B)=(ca_{ij})B=a_{ij}(cB)$.
\item For $A'=[a'_{ij}]$, $(a_{ij}+a'_{ij})B=a_{ij}B+a'_{ij}B$ and $a_{ij}(B+B')=a_{ij}B+a_{ij}B'$.
\item Write $C=[c_{\alpha\beta}]$. Under lexicographic indexing, the $((i,u,\alpha),(j,v,\beta))$ entry of each side is $a_{ij}b_{uv}c_{\alpha\beta}$.
\item The corresponding $(j,i)$ blocks are $(a_{ij}B)^t=a_{ij}B^t$.
\item The diagonal blocks of $I_m\otimes I_p$ are $I_p$, and its other blocks are zero; hence $I_m\otimes I_p=I_{mp}$.
\item Suppose $\mathbb{F}=\mathbb{C}$. The first identity holds by blockwise conjugation. By part (iv), $(A\otimes B)^*=(\overline{A}\otimes\overline{B})^t=\overline{A}^{\,t}\otimes\overline{B}^{\,t}=A^*\otimes B^*$.
\end{enumerate}
\hlt{Structural Commentary}: Entry factorization handles linearity, reindexing, and adjoints uniformly.\\
\hlt{Key Insight}: Each Kronecker entry is a product $a_{ij}b_{uv}$.
\end{proof}
\end{proposition}

\begin{theorem} \label{thm:kronecker-mixed-product-property} \hlt{Mixed Product Property}\\
Let $m,n,p,q,r,s\in\mathbb{Z}_{>0}$ and $A\in\operatorname{Mat}(m,n,\mathbb{F})$. Let $B\in\operatorname{Mat}(p,q,\mathbb{F})$ and $C\in\operatorname{Mat}(n,r,\mathbb{F})$.\\
Let $D\in\operatorname{Mat}(q,s,\mathbb{F})$. Then $(A\otimes B)(C\otimes D)=(AC)\otimes(BD)$.
\begin{proof}
Write $A=[a_{ik}]$ and $C=[c_{kj}]$, and fix $(i,j)\in J_m\times J_r$.\\
The $(i,j)$ block of $(A\otimes B)(C\otimes D)$ is $\sum\limits_{k=1}^{n}(a_{ik}B)(c_{kj}D)=\left(\sum\limits_{k=1}^{n}a_{ik}c_{kj}\right)BD=(AC)_{ij}BD$.\\
This is the corresponding block of $(AC)\otimes(BD)$, so the matrices are equal.\\
\hlt{Structural Commentary}: Block multiplication separates the contracted index from the factor product $BD$.\\
\hlt{Key Insight}: The $k$-sum is the $(i,j)$ entry of $AC$.
\end{proof}
\end{theorem}

\begin{corollary} \label{cor:inverse-of-kronecker-product} \hlt{Inverse of a Kronecker Product}\\
Let $m,n\in\mathbb{Z}_{>0}$, $A\in\operatorname{Mat}(m,m,\mathbb{F})^\times$, and $B\in\operatorname{Mat}(n,n,\mathbb{F})^\times$. Then $A\otimes B\in\operatorname{Mat}(mn,mn,\mathbb{F})^\times$ and $(A\otimes B)^{-1}=A^{-1}\otimes B^{-1}$.
\begin{proof}
Set $K\defeq A^{-1}\otimes B^{-1}$. The identity $(A\otimes B)K=(AA^{-1})\otimes(BB^{-1})=I_{mn}$ holds by Theorem \ref{thm:kronecker-mixed-product-property} and Proposition \ref{prop:basic-kronecker-identities}(v).\\
The identity $K(A\otimes B)=(A^{-1}A)\otimes(B^{-1}B)=I_{mn}$ holds by the same theorem and proposition.\\
\hlt{Structural Commentary}: The mixed-product law verifies both inverse equations factorwise.\\
\hlt{Key Insight}: The two-sided inverse is $A^{-1}\otimes B^{-1}$.
\end{proof}
\end{corollary}

\begin{proposition} \label{prop:kronecker-product-and-sum-diagonalizable-spectra} \hlt{Kronecker Product and Sum Spectra in the Diagonalizable Case}\\
Let $n,q\in\mathbb{Z}_{>0}$, $A\in\operatorname{Mat}(n,n,\mathbb{F})$, and $B\in\operatorname{Mat}(q,q,\mathbb{F})$. Let $P\in\operatorname{Mat}(n,n,\mathbb{F})^\times$ and $Q\in\operatorname{Mat}(q,q,\mathbb{F})^\times$.\\
Let $\lambda_1,\ldots,\lambda_n\in\mathbb{F}$ and $\mu_1,\ldots,\mu_q\in\mathbb{F}$. Set $\Lambda\defeq\operatorname{diag}(\lambda_1,\ldots,\lambda_n)$ and $M\defeq\operatorname{diag}(\mu_1,\ldots,\mu_q)$. Suppose $A=P\Lambda P^{-1}$ and $B=QMQ^{-1}$.
\begin{enumerate}[label=(\roman*)]
\setlength{\itemsep}{0pt}
\item $A\otimes B\sim\Lambda\otimes M$.
\item $A\otimes I_q+I_n\otimes B\sim\Lambda\otimes I_q+I_n\otimes M$.
\item $\operatorname{spec}(A\otimes B)=\{\lambda_i\mu_j:(i,j)\in J_n\times J_q\}$.
\item $\operatorname{spec}(A\otimes I_q+I_n\otimes B)=\{\lambda_i+\mu_j:(i,j)\in J_n\times J_q\}$.
\end{enumerate}
\begin{proof}
Set $R\defeq P\otimes Q$. Then $R\in\operatorname{Mat}(nq,nq,\mathbb{F})^\times$ and $R^{-1}=P^{-1}\otimes Q^{-1}$ by Corollary \ref{cor:inverse-of-kronecker-product}.
\begin{enumerate}[label=(\roman*)]
\setlength{\itemsep}{0pt}
\item The identity $A\otimes B=R(\Lambda\otimes M)R^{-1}$ holds by Theorem \ref{thm:kronecker-mixed-product-property}. Hence $A\otimes B\sim\Lambda\otimes M$.
\item By distributivity and Theorem \ref{thm:kronecker-mixed-product-property}, $A\otimes I_q+I_n\otimes B=R(\Lambda\otimes I_q+I_n\otimes M)R^{-1}$. Hence $A\otimes I_q+I_n\otimes B\sim\Lambda\otimes I_q+I_n\otimes M$.
\item The matrix $\Lambda\otimes M$ is diagonal with entries $(\lambda_i\mu_j)_{(i,j)\in J_n\times J_q}$.\\
Thus $\operatorname{spec}(A\otimes B)=\{\lambda_i\mu_j:(i,j)\in J_n\times J_q\}$ by part (i) and Theorem \ref{thm:similarity-preserves-characteristic-polynomial}(i).
\item The matrix $\Lambda\otimes I_q+I_n\otimes M$ is diagonal with entries $(\lambda_i+\mu_j)_{(i,j)\in J_n\times J_q}$.\\
Thus $\operatorname{spec}(A\otimes I_q+I_n\otimes B)=\{\lambda_i+\mu_j:(i,j)\in J_n\times J_q\}$ by part (ii) and Theorem \ref{thm:similarity-preserves-characteristic-polynomial}(i).
\end{enumerate}
\hlt{Structural Commentary}: The similarity $P\otimes Q$ transfers product and sum to diagonal coordinates.\\
\hlt{Key Insight}: Paired diagonal entries multiply for products and add for sums.
\end{proof}
\end{proposition}

For $r\in\mathbb{Z}_{>0}$, set $T_r\defeq\operatorname{tridiag}(-1,2,-1)\in\operatorname{Mat}(r,r,\mathbb{R})$.\\
For $(m,n)\in\mathbb{Z}_{>0}^2$, Figure \ref{fig:kronecker-sum-five-point-stencil} depicts the interior stencil of $T_m\otimes I_n+I_m\otimes T_n$.\\
The two summands act in the vertical and horizontal coordinate directions, respectively.

\begin{illustrationfigure}[H]
  \centering
  \begin{tikzpicture}[
  >=Stealth,
  every node/.style={font=\small},
  gridpt/.style={circle, draw=gray!35, fill=gray!3, minimum size=0.32cm, inner sep=0pt},
  vpt/.style={circle, draw=blue!65!black, fill=blue!7, minimum size=0.48cm, inner sep=0pt, font=\scriptsize},
  hpt/.style={circle, draw=teal!65!black, fill=teal!8, minimum size=0.48cm, inner sep=0pt, font=\scriptsize},
  sumpt/.style={circle, draw=blue!65!black, fill=blue!7, minimum size=0.48cm, inner sep=0pt, font=\scriptsize},
  link/.style={gray!35, line width=0.42pt},
  vlink/.style={blue!70!black, line width=0.85pt},
  hlink/.style={teal!70!black, line width=0.85pt}
]
  \newcommand{\plainstencil}[2]{%
    \foreach \x in {-0.55,0,0.55} {
      \foreach \y in {-0.55,0,0.55} {
        \node[gridpt] at (#1+\x,#2+\y) {};
      }
    }
    \draw[link] (#1-0.55,#2-0.55) -- (#1-0.55,#2+0.55);
    \draw[link] (#1,#2-0.55) -- (#1,#2+0.55);
    \draw[link] (#1+0.55,#2-0.55) -- (#1+0.55,#2+0.55);
    \draw[link] (#1-0.55,#2-0.55) -- (#1+0.55,#2-0.55);
    \draw[link] (#1-0.55,#2) -- (#1+0.55,#2);
    \draw[link] (#1-0.55,#2+0.55) -- (#1+0.55,#2+0.55);
  }

  \plainstencil{0}{0}
  \draw[vlink] (0,-0.55) -- (0,0.55);
  \node[vpt] at (0,0.55) {$-1$};
  \node[vpt] at (0,0) {$2$};
  \node[vpt] at (0,-0.55) {$-1$};
  \node at (0,-1.08) {$T_m\otimes I_n$};

  \node at (1.65,0) {$+$};

  \plainstencil{3.3}{0}
  \draw[hlink] (2.75,0) -- (3.85,0);
  \node[hpt] at (2.75,0) {$-1$};
  \node[hpt] at (3.3,0) {$2$};
  \node[hpt] at (3.85,0) {$-1$};
  \node at (3.3,-1.08) {$I_m\otimes T_n$};

  \node at (4.95,0) {$=$};

  \plainstencil{6.6}{0}
  \draw[vlink] (6.6,-0.55) -- (6.6,0.55);
  \draw[hlink] (6.05,0) -- (7.15,0);
  \node[sumpt] at (6.6,0.55) {$-1$};
  \node[sumpt] at (6.05,0) {$-1$};
  \node[sumpt] at (6.6,0) {$4$};
  \node[sumpt] at (7.15,0) {$-1$};
  \node[sumpt] at (6.6,-0.55) {$-1$};
  \node at (6.6,-1.08) {five-point stencil};
\end{tikzpicture}

  \caption{The two tridiagonal factors combine into center weight $4$ and four neighboring weights $-1$.}
  \label{fig:kronecker-sum-five-point-stencil}
\end{illustrationfigure}

\begin{definition} \label{def:vectorization-map} \hlt{Vectorisation}\\
Let $m,n\in\mathbb{Z}_{>0}$ and $X=[\mathbf{x}_1\ \mathbf{x}_2\ \cdots\ \mathbf{x}_n]\in\operatorname{Mat}(m,n,\mathbb{F})$. The \hlt{vectorisation map} is $\operatorname{vec}\colon\operatorname{Mat}(m,n,\mathbb{F})\to\mathbb{F}^{mn}$.\\
For $X$ as above, $\operatorname{vec}(X)\defeq(\mathbf{x}_1^t\ \cdots\ \mathbf{x}_n^t)^t$.
\end{definition}

\begin{theorem} \label{thm:vectorization-of-triple-product} \hlt{Vectorisation of a Triple Product}\\
Let $m,n,p,q\in\mathbb{Z}_{>0}$ and $A\in\operatorname{Mat}(m,n,\mathbb{F})$. Let $X\in\operatorname{Mat}(n,p,\mathbb{F})$ and $B\in\operatorname{Mat}(p,q,\mathbb{F})$. Then $\operatorname{vec}(AXB)=(B^t\otimes A)\operatorname{vec}(X)$.
\begin{proof}
Write $B=[b_{ik}]=[\mathbf{b}_1\ \cdots\ \mathbf{b}_q]$ and $X=[\mathbf{x}_1\ \cdots\ \mathbf{x}_p]$.\\
For $k\in J_q$, the $k$th block on each side is $AX\mathbf{b}_k=\sum\limits_{i=1}^{p}b_{ik}A\mathbf{x}_i$; hence the vectors are equal.\\
\hlt{Structural Commentary}: Column stacking matches the $k$th block with the $k$th column of $AXB$.\\
\hlt{Key Insight}: Both blocks equal $\sum\limits_{i=1}^{p}b_{ik}A\mathbf{x}_i$.
\end{proof}
\end{theorem}

\begin{perspective} \label{per:vectorized-minimal-polynomial-linear-system} \hlt{Vectorising a Minimal-Polynomial Test}\\
The candidate step in Method \ref{meth:operator-minimal-polynomial-from-matrix} admits a vectorised form. Fix $n\in\mathbb{Z}_{>0}$, $A\in\operatorname{Mat}(n,n,\mathbb{F})$, and $m\in\{1,\ldots,n^2\}$.\\
For $(a_0,\ldots,a_{m-1})\in\mathbb{F}^m$, set $p_m(z)\defeq a_0+a_1z+\cdots+a_{m-1}z^{m-1}+z^m$. Set $V_m(A)\defeq[\operatorname{vec}(I_n)\ \operatorname{vec}(A)\ \cdots\ \operatorname{vec}(A^{m-1})]\in\operatorname{Mat}(n^2,m,\mathbb{F})$.\\
Set $\mathbf{a}\defeq(a_0,\ldots,a_{m-1})^t$. Then $p_m(A)=0\iff V_m(A)\mathbf{a}=-\operatorname{vec}(A^m)$.\\
\hlt{Structural Commentary}: Vectorisation turns the candidate equation into an $n^2\times m$ scalar system.\\
\hlt{Key Insight}: Its columns are vectorised matrix powers.
\end{perspective}
